\paragrafos{Θέσεις ευθείας και κύκλου: εξωτερική, εφαπτομένη, τέμνουσα}{Β.1.13, «Θέσεις ευθείας και κύκλου»}

\begin{anagke}
Γύρω από το στρογγυλό παρτέρι θα περάσει ένα μονοπάτι, και η Α3 συζητάει πού θα το τοποθετήσει. Ο Νίκος προτείνει να περνά ακριβώς πάνω από το παρτέρι. «Θα πατάμε τα λουλούδια!», διαμαρτύρεται η Ρήνη. Ο Βασίλης προτείνει ένα μονοπάτι που περνά μακριά από το παρτέρι, χωρίς να το αγγίζει. «Τότε δεν θα το βλέπουμε από κοντά», λέει η Αγγελική. Η Πολυξένη προτείνει το τρίτο: ένα μονοπάτι που ακουμπά το παρτέρι σε ένα μόνο σημείο, σαν τον δρόμο που ακουμπά μια ρόδα.

Η κυρία Αλεξάνδρα το ονομάζει: τρεις θέσεις μιας ευθείας ως προς έναν κύκλο. Δείχνει πώς καταλαβαίνουμε ποια είναι από την απόσταση του κέντρου από την ευθεία: «Μετρήστε την απόσταση του κέντρου από το μονοπάτι και συγκρίνετέ την με την ακτίνα».

Στο τέλος του κεφαλαίου δένουμε την απόσταση σημείου από ευθεία (1.10) με τον κύκλο (1.11): τρεις θέσεις ευθείας και κύκλου, και τι ισχύει όταν η ευθεία εφάπτεται.
\end{anagke}

\begin{lexi}{εφαπτομένη, τέμνουσα}{εφαπτομένη, τέμνουσα, εξωτερική}
Η \lexiword{εφαπτομένη} είναι αυτή που «εφάπτεται», δηλαδή «ακουμπά», μόνο σε ένα σημείο. Η \lexiword{τέμνουσα} είναι αυτή που «τέμνει», περνά μέσα από τον κύκλο και τον τέμνει σε δύο σημεία. Η \lexiword{εξωτερική} ευθεία βρίσκεται «έξω» από τον κύκλο και δεν τον ακουμπά.
\end{lexi}

\begin{metafrasi}
Η απόσταση του κέντρου από την ευθεία. & $\mathrm{ΟΜ}=d$\\
Η ευθεία είναι μακριά από τον κύκλο. & $d>\rho$: εξωτερική\\
Η ευθεία ακουμπά τον κύκλο σε ένα σημείο. & $d=\rho$: εφαπτομένη\\
Η ευθεία περνά μέσα στον κύκλο. & $d<\rho$: τέμνουσα\\
\end{metafrasi}

\tmima{Η θεωρία}

Έστω κύκλος $(\mathrm{Ο},\rho)$ και ευθεία $(\varepsilon)$ και $d=\mathrm{ΟΜ}$ η απόσταση του κέντρου από την ευθεία (με $\mathrm{ΟΜ}\perp\varepsilon$). Υπάρχουν τρεις περιπτώσεις:

\begin{center}
\begin{tikzpicture}[scale=0.85]
\begin{scope}
\draw (0,0) circle (1.2); \fill (0,0) circle (1.5pt) node[above left,font=\small] {$\mathrm{Ο}$};
\draw[thick,{Stealth}-{Stealth}] (-1.8,1.9)--(1.8,1.9); \draw[gblue] (0,0)--(0,1.9);
\node at (0,-1.7) {\small $d>\rho$: εξωτερική};
\end{scope}
\begin{scope}[xshift=4.5cm]
\draw (0,0) circle (1.2); \fill (0,0) circle (1.5pt);
\draw[thick,{Stealth}-{Stealth}] (-1.8,1.2)--(1.8,1.2); \draw[gblue] (0,0)--(0,1.2); \fill (0,1.2) circle (2pt) node[above right,font=\small] {$\mathrm{Μ}$};
\node at (0,-1.7) {\small $d=\rho$: εφαπτομένη};
\end{scope}
\begin{scope}[xshift=9cm]
\draw (0,0) circle (1.2); \fill (0,0) circle (1.5pt);
\draw[thick,{Stealth}-{Stealth}] (-1.8,0.6)--(1.8,0.6); \draw[gblue] (0,0)--(0,0.6);
\fill (-1.04,0.6) circle (2pt) node[above left,font=\small] {$\mathrm{Α}$}; \fill (1.04,0.6) circle (2pt) node[above right,font=\small] {$\mathrm{Β}$};
\node at (0,-1.7) {\small $d<\rho$: τέμνουσα};
\end{scope}
\end{tikzpicture}
\end{center}

\begin{itemize}[leftmargin=1.2em,itemsep=0.3ex,topsep=2pt]
\item Αν $d>\rho$, η ευθεία δεν έχει κανένα κοινό σημείο με τον κύκλο: είναι \textbf{εξωτερική} του κύκλου.
\item Αν $d=\rho$, η ευθεία έχει \textbf{ένα} κοινό σημείο $\mathrm{Μ}$ με τον κύκλο και λέγεται \textbf{εφαπτομένη} του κύκλου στο $\mathrm{Μ}$. Το $\mathrm{Μ}$ λέγεται σημείο επαφής και η εφαπτομένη είναι \textbf{κάθετη} στην ακτίνα $\mathrm{ΟΜ}$.
\item Αν $d<\rho$, η ευθεία έχει \textbf{δύο} κοινά σημεία $\mathrm{Α}$, $\mathrm{Β}$ με τον κύκλο και λέγεται \textbf{τέμνουσα}.
\end{itemize}

Αντίστροφα, από το πλήθος των κοινών σημείων προκύπτει η θέση: $0$ κοινά σημεία σημαίνει εξωτερική, $1$ εφαπτομένη, $2$ τέμνουσα. Μια ευθεία δεν μπορεί να έχει με κύκλο περισσότερα από δύο κοινά σημεία.

\begin{prosoxi}
\begin{itemize}[leftmargin=1.2em,itemsep=0.3ex,topsep=2pt]
\item Η απόσταση του κέντρου από την ευθεία είναι η κάθετη απόσταση, όχι οποιοδήποτε πλάγιο τμήμα.
\item Η εφαπτομένη έχει ένα μόνο κοινό σημείο με τον κύκλο. Αν έχει δύο, είναι τέμνουσα.
\item Η εφαπτομένη είναι κάθετη στην ακτίνα του σημείου επαφής, όχι σε οποιαδήποτε ακτίνα.
\item Σε ένα σημείο του κύκλου υπάρχει μία και μοναδική εφαπτομένη.
\end{itemize}
\end{prosoxi}

\tmima{Λυμένα παραδείγματα}

\paradeigma{Η θέση του μονοπατιού}
\tixero Κύκλος $(\mathrm{Ο},3)$ και ευθείες με αποστάσεις $d$ από το κέντρο: $5$, $3$, $1{,}5$.\par
\tithelo Η θέση καθεμίας.\par
\possindeo Συγκρίνω το $d$ με το $\rho=3$. $5>3$: εξωτερική. $3=3$: εφαπτομένη. $1{,}5<3$: τέμνουσα.\par
\apantisi Εξωτερική, εφαπτομένη, τέμνουσα.\par
\elegxos Στην τέμνουσα υπάρχουν δύο κοινά σημεία.

\paradeigma{Κοινά σημεία}
\tixero Μια ευθεία με απόσταση $d=4$ από το κέντρο κύκλου ακτίνας $\rho=4$.\par
\tithelo Πόσα κοινά σημεία έχει.\par
\possindeo $d=\rho$, εφαπτομένη, άρα ένα κοινό σημείο.\par
\apantisi Ένα.\par
\elegxos Το σημείο επαφής είναι το ίχνος της καθέτου από το κέντρο.

\paradeigma{Η εφαπτομένη σε σημείο}
\tixero Ένα σημείο $\mathrm{Μ}$ του κύκλου $(\mathrm{Ο},\rho)$.\par
\tithelo Πώς σχεδιάζω την εφαπτομένη στο $\mathrm{Μ}$;\par
\possindeo Φέρνω την ακτίνα $\mathrm{ΟΜ}$ και την κάθετη στην $\mathrm{ΟΜ}$ στο σημείο $\mathrm{Μ}$.\par
\apantisi Η εφαπτομένη είναι η κάθετη στην ακτίνα $\mathrm{ΟΜ}$ στο $\mathrm{Μ}$.\par
\elegxos Η απόσταση του $\mathrm{Ο}$ από την ευθεία είναι $\mathrm{ΟΜ}=\rho$.

\paradeigma{Πόση απόσταση;}
\tixero Κύκλος $(\mathrm{Ο},5)$ και ευθεία που τον τέμνει σε δύο σημεία.\par
\tithelo Τι ξέρουμε για την απόσταση $d$ του κέντρου από την ευθεία.\par
\possindeo Αφού είναι τέμνουσα, $d<5$.\par
\apantisi $0\le d<5$.\par
\elegxos Αν $d=0$ η ευθεία περνά από το κέντρο και είναι διάμετρος.

\tmima{Ασκήσεις}

\zoni{Το κατάλαβα}{τρεις θέσεις}

\ex Κύκλος $(\mathrm{Ο},3)$. Πες τη θέση της ευθείας όταν η απόσταση του κέντρου από αυτήν είναι: α) $5$ \quad β) $3$ \quad γ) $1{,}5$ \quad δ) $0$ \quad ε) $3{,}01$.

\ex Πόσα κοινά σημεία έχει ευθεία με κύκλο όταν είναι εξωτερική, εφαπτομένη, τέμνουσα;

\ex Σχεδίασε κύκλο με ακτίνα $3$ cm και φέρε μια εξωτερική, μια εφαπτομένη και μια τέμνουσα. Μέτρησε την απόσταση του κέντρου από την καθεμία.

\ex Σχεδίασε την εφαπτομένη ενός κύκλου σε ένα σημείο του $\mathrm{Μ}$.

\ex Κύκλος $(\mathrm{Ο},5)$ και ευθεία που τον τέμνει σε δύο σημεία. Μεταξύ ποιων τιμών βρίσκεται η απόσταση $d$ του κέντρου από την ευθεία;

\ex Σωστό ή λάθος; α) Η εφαπτομένη έχει ένα κοινό σημείο με τον κύκλο. \quad β) Η τέμνουσα έχει τρία κοινά σημεία. \quad γ) Η εφαπτομένη είναι κάθετη στην ακτίνα του σημείου επαφής. \quad δ) Αν $d>\rho$ η ευθεία τέμνει τον κύκλο.

\zoni{Το χρησιμοποιώ}{μονοπάτια και ρόδες}

\ex \textbf{Το μονοπάτι.} Το παρτέρι έχει ακτίνα $2{,}5$ μ. Ένα ευθύγραμμο μονοπάτι περνά σε απόσταση $2{,}5$ μ. από το κέντρο. Πατάμε τα λουλούδια;

\ex \textbf{Τρία μονοπάτια.} Τρία ευθύγραμμα μονοπάτια περνούν σε αποστάσεις $1{,}8$ μ., $2{,}5$ μ. και $3{,}4$ μ. από το κέντρο του παρτεριού ακτίνας $2{,}5$ μ. Ποιο τέμνει το παρτέρι, ποιο το εφάπτεται και ποιο δεν το ακουμπά;

\ex \textbf{Ο τροχός.} Ένας τροχός ακτίνας $33$ cm κυλά σε ευθύγραμμο δρόμο. Πόσο απέχει το κέντρο του από το δρόμο; Πόσα κοινά σημεία έχει ο τροχός με τον δρόμο;

\ex \textbf{Ο φακός.} Ένας φακός φωτίζει ευθεία γραμμή που εφάπτεται σε ένα κυκλικό δωμάτιο. Πού ακουμπά το φως στον τοίχο;

\ex \textbf{Η ρόδα.} Γιατί μια ρόδα κυλά ομαλά πάνω σε επίπεδο δρόμο; (σκέψου ότι ο δρόμος εφάπτεται στη ρόδα).

\zoni{Το σκέφτομαι}{όταν η ευθεία ακουμπά}

\ex Γιατί η εφαπτομένη είναι κάθετη στην ακτίνα του σημείου επαφής; (Σκέψου ότι είναι η μικρότερη απόσταση του κέντρου από την ευθεία.)

\ex Ευθεία περνά από δύο διαφορετικά σημεία του κύκλου. Μπορεί να έχει και τρίτο κοινό σημείο; Εξήγησε.

\ex Πόσες εφαπτομένες διέρχονται από σημείο του κύκλου; Και από σημείο εκτός κύκλου; (Σχεδίασε.)

\ex Δύο παράλληλες ευθείες εφάπτονται σε κύκλο ακτίνας $\rho$. Πόση είναι η απόσταση ανάμεσά τους;

\zoni{Για το σπίτι}{εφαπτόμενα σχήματα}

\ex Βρες στο σπίτι πέντε περιπτώσεις ευθείας και κύκλου με κοινό ένα σημείο (π.χ. τροχός πάνω στο πάτωμα).

\ex Σχεδίασε έναν κύκλο και τρεις εφαπτομένες του σε διαφορετικά σημεία.
